% Main-text Table 1: movement components required by representative scales.
% Required packages: booktabs, array, ragged2e, graphicx (for \rotatebox).

\begin{table}[!htbp]
\centering
\caption{Score-defining movement components of representative post-stroke scales. A component is marked only when the scoring criteria grade it; shared components do not make scores interchangeable.}
\label{tab:representative_scales_compact}
\begingroup
\newcommand{\mcmark}{\(\bullet\)}
\setlength{\tabcolsep}{4pt}
\renewcommand{\arraystretch}{1.1}
\small
\begin{tabular}{@{}l l l *{7}{c}@{}}
\toprule
 & & & \multicolumn{6}{c}{\textbf{Movement component}} & \textbf{Modifier} \\
\cmidrule(lr){4-9}\cmidrule(l){10-10}
\textbf{Scale} & \textbf{COA} & \textbf{Score output} & ROM & S/F & Coord. & Bal. & Dex. & Gait & Q/C \\
\midrule
\multicolumn{10}{@{}l}{\textit{Upper limb}} \\
FMA-UE & ClinRO & Ordinal (0--66) & \mcmark & \mcmark & \mcmark &  & \mcmark &  & \mcmark \\
ARAT & ClinRO & Ordinal (0--57) &  &  & \mcmark &  & \mcmark &  & \mcmark \\
\midrule
\multicolumn{10}{@{}l}{\textit{Trunk and balance}} \\
TIS & ClinRO & Ordinal (0--23) &  &  & \mcmark & \mcmark &  &  & \mcmark \\
BBS & ClinRO & Ordinal (0--56) &  &  &  & \mcmark &  & \mcmark & \mcmark \\
Mini-BESTest & ClinRO & Ordinal (0--28) & \mcmark &  &  & \mcmark &  & \mcmark & \mcmark \\
\midrule
\multicolumn{10}{@{}l}{\textit{Lower limb and mobility}} \\
FMA-LE & ClinRO & Ordinal (0--34) & \mcmark & \mcmark & \mcmark &  &  &  & \mcmark \\
DGI & ClinRO & Ordinal (0--24) &  &  & \mcmark & \mcmark &  & \mcmark & \mcmark \\
FAC & ClinRO & Category (0--5) &  &  &  & \mcmark &  & \mcmark &  \\
10MWT & PerfO & Speed (m/s) &  &  &  &  &  & \mcmark &  \\
6MWT & PerfO & Distance (m) &  &  &  &  &  & \mcmark &  \\
TUG\textsuperscript{a} & PerfO & Time (s) &  &  &  &  &  & \mcmark &  \\
\midrule
\multicolumn{10}{@{}l}{\textit{Global outcome (movement-related items only)}} \\
NIHSS & ClinRO & Ordinal (0--42) &  & \mcmark & \mcmark &  &  &  &  \\
BI\textsuperscript{a} & ClinRO & Ordinal (0--100) &  &  &  &  &  & \mcmark &  \\
\bottomrule
\end{tabular}

\vspace{3pt}
\begin{minipage}{\linewidth}
\footnotesize
\textit{Notes.} \(\bullet\), at least one item's criteria grade performance in that component (rules in Supplementary Table~\ref{tab:supplementary_tool_library_motion}); components a task requires without changing the grade are not marked, such as balance in the time-scored 10MWT and TUG or trunk range in TIS items graded by completion, symmetry and compensation. Gait includes transfers and stepping (BBS, BI). S/F covers FMA items graded against examiner resistance and, for the NIHSS, holding a limb against gravity. Q/C, movement quality and compensation, is a cross-component modifier marked only when criteria grade quality, abnormal synergy or compensation. NIHSS: motor and limb-ataxia items; BI: mobility items. \textsuperscript{a}Guideline-supported tools added to the 11 SRRR-priority tools for transitional mobility (TUG) and daily activities (BI). COA is an operational classification under FDA--NIH definitions (Section~\ref{sec:scales}) \cite{FDANIH2016BESTCOA,Walton2015COAConceptualFoundation}; under a task-based reading, the task-rated ClinRO tools here (FMA-UE, ARAT, TIS, BBS, Mini-BESTest, FMA-LE, DGI) would be PerfO (Supplementary Table~\ref{tab:supplementary_tool_library_clinical}); components were informed by the ICF, guidelines and SRRR domains \cite{WHO2001ICF,S1Winstein2016AHAASARehab,S1VADoD2024StrokeRehab,S1NICE2023NG236,CoreKwakkel2017SRRRSensorimotor,CoreKwakkel2019SRRRUpperLimb,CoreVanCriekinge2024SRRRBalanceMobility}.
ROM, joint mobility and range of motion; S/F, strength and force production; Coord., coordination and selective motor control; Bal., balance and postural control; Dex., manual dexterity and upper-limb task execution; Gait, gait and functional mobility.
10MWT, 10-Metre Walk Test; 6MWT, 6-Minute Walk Test; ARAT, Action Research Arm Test; BBS, Berg Balance Scale; BI, Barthel Index; ClinRO, clinician-reported outcome; COA, clinical outcome assessment; DGI, Dynamic Gait Index; FAC, Functional Ambulation Category; FMA-LE/UE, Fugl--Meyer Assessment lower/upper extremity; Mini-BESTest, Mini-Balance Evaluation Systems Test; NIHSS, National Institutes of Health Stroke Scale; PerfO, performance outcome; TIS, Trunk Impairment Scale; TUG, Timed Up and Go.
\end{minipage}
\endgroup
\end{table}
